% Ad M. Brutum 1.16 --- headnote
% ad-brutum-01-16, early May 43 BC, in the East

\textit{M. Brutus to Cicero himself --- the salutation reads
\textit{Brutus Ciceroni salutem} --- written from the East, where
Brutus held his eastern command, at the beginning of May 43 BC. The
Perseus dateline reads \textit{Scr. Athenis in. Maio a. 711 (43)},
``written at Athens, at the beginning of May,'' 43. The meta entry's
year-precision placeholder is here corrected to early May 43; the
letter belongs to the spring of that year, months before the autumn
proscriptions, and the place-note ``at Athens'' is the conventional
ascription rather than a fixed fact.}

\textit{The occasion is a letter Cicero had written to the young
Caesar --- Octavian, Caesar's heir --- a fragment of which Atticus
had forwarded to Brutus. In it Cicero, thanking Octavian for his
services to the state, had commended Brutus's own safety to the
boy's care, in terms Brutus found suppliant and abject. The whole
of this letter is Brutus's answer, and it is among the most severe
documents in the corpus. Its argument is single and relentless: that
to beg the safety of the liberators of the world from one man is to
confess that the despotism has not been abolished but only its master
changed (\textit{non sublatam dominationem sed dominum commutatum
esse}); that Cicero's words, looked at honestly, are the prayers of
a slave addressed to a king (\textit{servientis adversus regem
preces}); and that safety bought at the price of dignity and liberty
is a death worse than death. Brutus refuses it for himself --- he
will hold Rome to be wherever a man may be free --- and warns Cicero
that the courage shown against Antony will curdle into a reputation
for cowardice if it is now bent toward courting Octavian as a more
agreeable master.}

\textit{The letter is the companion to \textit{Ad M. Brutum} 1.17,
Brutus's contemporary letter to Atticus, which makes the same case
in private and at greater heat; the two survived together in the
same archive and were transmitted as a pair. The text carries an
unresolved crux among the corrupt loci of the central sections, here
translated through by sense without brackets. Within a year and a
half Brutus's prediction came true: Cicero fell under the
proscriptions of the master he had helped to make.}
